\begin{tikzpicture}[
  font=\scriptsize,
  execute at begin node={\hyphenpenalty=10000\exhyphenpenalty=10000},
  tahap/.style={draw, rounded corners, fill=blue!8, text width=1.7cm, minimum height=0.95cm, align=center},
  tekan/.style={tahap, fill=orange!12, text width=1.5cm, minimum height=2.0cm},
  hasil/.style={tahap, fill=green!10},
  ket/.style={font=\tiny, align=center},
  panah/.style={-{Stealth[length=2mm]}, thick},
  putus/.style={{Stealth[length=2mm]}-{Stealth[length=2mm]}, thin, dashed}
]
  % Lingkungan layanan digambar lebih dulu supaya berada di belakang node.
  \draw[rounded corners, fill=praditagreen!15] (4.6,1.85) rectangle (11.75,-1.85);
  \node[ket] at (8.175,1.6) {\textbf{Lingkungan layanan Medikara: komunikasi antarlayanan}};

  \node[tahap] (klinik) at (0,0.75) {40 klinik mitra};
  \node[tahap] (plus) at (0,-0.75) {Medikara+};
  \node[tekan] (gerbang) at (3.45,0) {\textit{API gateway}\\{\tiny\mdseries batas akses API}};

  \draw[panah] (klinik.east) -- node[above, ket] {lalu lintas masuk} (gerbang.west |- klinik);
  \draw[panah] (plus.east) -- node[below, ket] {token, laju, kuota} (gerbang.west |- plus);

  \node[tahap] (janji) at (5.85,0.65) {Layanan janji temu};
  \node[tahap] (hasillab) at (9.35,0.65) {Layanan hasil};
  \node[tahap] (identitas) at (5.85,-0.95) {Layanan identitas};
  \node[hasil] (adaptor) at (9.35,-0.95) {Adaptor SIMRS};

  \draw[panah] (gerbang.east |- janji) -- (janji);
  \draw[panah] (gerbang.east |- identitas) -- (identitas);

  \draw[putus] (janji) -- node[above, ket] {kendali\\per layanan} (hasillab);
  \draw[putus] (janji) -- (identitas);
  \draw[putus] (janji) -- (adaptor);
  \draw[putus] (hasillab) -- (adaptor);
  \draw[putus] (identitas) -- node[below, ket] {atau\\\textit{service mesh}} (adaptor);

  \node[ket, text width=1.0cm] at (11.15,-0.15) {rinci di IT40713};
\end{tikzpicture}
